% !TEX TS-program = pdflatex
% !TEX encoding = UTF-8 Unicode

% This is a simple template for a LaTeX document using the "article" class.
% See "book", "report", "letter" for other types of document.

\documentclass[11pt,titlepage]{article} % use larger type; default would be 10pt

\usepackage[utf8]{inputenc} % set input encoding (not needed with XeLaTeX)

%%% Examples of Article customizations
% These packages are optional, depending whether you want the features they provide.
% See the LaTeX Companion or other references for full information.

%%% PAGE DIMENSIONS
\usepackage{geometry} % to change the page dimensions
\geometry{a4paper} % or letterpaper (US) or a5paper or....
% \geometry{margin=2in} % for example, change the margins to 2 inches all round
% \geometry{landscape} % set up the page for landscape
%   read geometry.pdf for detailed page layout information

\usepackage{graphicx} % support the \includegraphics command and options
\usepackage{titlepic}

% \usepackage[parfill]{parskip} % Activate to begin paragraphs with an empty line rather than an indent

%%% PACKAGES
\usepackage{booktabs} % for much better looking tables
\usepackage{array} % for better arrays (eg matrices) in maths
\usepackage{paralist} % very flexible & customisable lists (eg. enumerate/itemize, etc.)
\usepackage{verbatim} % adds environment for commenting out blocks of text & for better verbatim
\usepackage{subfig} % make it possible to include more than one captioned figure/table in a single float
% These packages are all incorporated in the memoir class to one degree or another...

%%% HEADERS & FOOTERS
\usepackage{fancyhdr} % This should be set AFTER setting up the page geometry
\pagestyle{plain} % options: empty , plain , fancy
\renewcommand{\headrulewidth}{0pt} % customise the layout...
\lhead{}\chead{}\rhead{}
\lfoot{}\cfoot{\thepage}\rfoot{}

%%% SECTION TITLE APPEARANCE
\usepackage{sectsty}
\allsectionsfont{\sffamily\mdseries\upshape} % (See the fntguide.pdf for font help)
% (This matches ConTeXt defaults)

%%% ToC (table of contents) APPEARANCE
\usepackage[nottoc,notlof,notlot]{tocbibind} % Put the bibliography in the ToC
\usepackage[titles,subfigure]{tocloft} % Alter the style of the Table of Contents
\renewcommand{\cftsecfont}{\rmfamily\mdseries\upshape}
\renewcommand{\cftsecpagefont}{\rmfamily\mdseries\upshape} % No bold!

\newenvironment{changemargin}[3]{%
\begin{list}{}{%
\setlength{\topsep}{0pt}%
\setlength{\headsep}{#3}%
\setlength{\leftmargin}{#1}%
\setlength{\rightmargin}{#2}%
\setlength{\listparindent}{\parindent}%
\setlength{\itemindent}{\parindent}%
\setlength{\parsep}{\parskip}%
}%
\item[]}{\end{list}}

%Table Formatting
\usepackage{tabularx,hhline}
\usepackage{pbox}

\usepackage{listings}
\usepackage{color}

\definecolor{myblue}{rgb}{0.031,0.56,0.741}
\definecolor{mygray}{rgb}{0.5,0.5,0.5}
\definecolor{mymauve}{rgb}{0.58,0,0.82}
\definecolor{lightgray}{rgb}{0.988,0.988,0.988}

\lstset{
   backgroundcolor=\color{lightgray},
   extendedchars=true,
   basicstyle=\footnotesize\ttfamily,
   showstringspaces=false,
   showspaces=false,
   numbers=left,
   numberstyle=\footnotesize,
   numbersep=9pt,
   tabsize=2,
   breaklines=true,
   showtabs=false,
   captionpos=b,
   commentstyle=\color{myblue},
   keywordstyle=\color{blue},
   stringstyle=\color{mymauve},
}

%%% END Article customizations

%%% The "real" document content comes below...

\titlepic{\includegraphics[scale=0.60]{polimi_logo.jpg}}
\title{Inspection \\ \vspace{1cm} \large{Version 1.0}} 
\author{Giorgio Pea(Mat. 853872), Andrea Sessa(Mat. 850082)}
\date{5/1/2016} 

\begin{document}

\maketitle

\newpage

\tableofcontents

\newpage

\section{Introduction}

\section{Classes}

Included in this section the two java classes subjected to the analisys.\newline

\noindent File: /appserver/web/web-core/src/main/java/org/apache/catalina/ssi/\textbf{SSIServlet.java}\newline\newline
Methods under inspection:
  \begin{itemize}
    \item \textbf{\textit{init()}}
    \item \textbf{\textit{requestHandler(HttpServletRequest req, HttpServletResponse res)}}
    \item \textbf{\textit{processSSI( HttpServletRequest req , HttpServletResponse res , URL resource )}}
  \end{itemize}

\noindent File: /appserver/web/web-core/src/main/java/org/apache/catalina/ssi/\textbf{SSIMediator.java}\newline\newline
Methods under inspection:
  \begin{itemize}
    \item \textbf{\textit{substituteVariables(String val)}}
  \end{itemize}
  
\newpage


\section{Functional Role}
  In this section are included some information about the functioning of tha analyzed classes and methods.\newline

\subsection{SSIServlet.java}

  From the Javadoc:
  
  \lstinputlisting[language=Java, firstnumber=78, firstline=78, lastline=81]{SSIServlet.java}

  \noindent This class represents a Java EE servlet used to process requests that include some SSI instruction.\newline
  SSI(Server Side Include) that is a simple interpreted server-side scripting language.
  The most frequent use of SSI is to include the contents of one or more files into a web page on a web server.

  \begin{itemize}
    \item \textit{Init()}
    \item \textit{requestHandler()}
      \newline From the inspection of the code this function is only called when the servelet receives a http Get or Post request.
      The javadoc for the method, included in the code, states:
      
      \lstinputlisting[language=Java, firstnumber=173, firstline=173, lastline=180]{SSIServlet.java}
      
      Hence the method accepts as parameters a HttpServletRequest,the incoming request, and a HttpServletResponse that is a reference
      to the response.\newline
      
      Now the objective of the method is to retrieve the correct resource from the servelet context.
      If the debug level is greater than zero then log a message into the logger for debug purposes.
      \lstinputlisting[language=Java, firstnumber=183, firstline=183, lastline=188]{SSIServlet.java}

      The comment is very clear: it checks if the resource is either in the `WEB-INF` or `META-INF` subdirectories;
      if so the function return with an error code.
      \lstinputlisting[language=Java, firstnumber=189, firstline=189, lastline=196]{SSIServlet.java}
      
      Here the fucntion tries to retrieve the URL to the resource; it also performs an existence check on the resource,
      if the resource doesn't exist the function return an error.
      \lstinputlisting[language=Java, firstnumber=197, firstline=197, lastline=202]{SSIServlet.java}
      
      In the final part, the function starts to initialize the header of the HttpServletResponse by setting:
      the mime type, the encoding of the output text and the expiration time for the response(in seconds, see init()).\newline
      Finally the processSSI() function is invoked passing as parameters the original request, the reference to the response and the resource.
      \lstinputlisting[language=Java, firstnumber=203, firstline=203, lastline=213]{SSIServlet.java}
      \newpage
    \item \textit{processSSI()}

  \end{itemize}

\subsection{SSIMediator.java}



\section{Issues}

\section{Additional Considerations}

\end{document}
